\section{总结与延伸}

\begin{importantbox}{本场核心观点速览}
\begin{enumerate}
  \item \textbf{架构}：MoE 是短期工程解，Dense 是长期方向；FLOPs 永远最便宜。
  \item \textbf{多模态}：Native Multimodal 长期看好，LLM/VLM 分离是暂时状态。输入端统一无争议，输出端仍有分歧。
  \item \textbf{数据}：合成数据面临跨代迁移和 Model Collapse 风险；Mid-training 是 World Model 建模的训练中间态。Vision-Centric 自监督可能破解多模态数据瓶颈。
  \item \textbf{RL}：闭环控制特性使得问题定义比算法更重要；缺乏 Scaling Law、泛化性不足、environment 复杂度是三大挑战。RL 正走向专业化。
  \item \textbf{Agent}：决策质量低、Memory 管理是核心瓶颈；Claude Code 是 2025 标杆；直接推向真实世界获取 reward signal 可能是突破口。
  \item \textbf{持续学习}：Context-based（Prompt Database + Retrieval）比 Parameter-based 更安全实用；长 context 本质还是 FLOPs。
\end{enumerate}
\end{importantbox}

\subsection{拓展阅读}
\begin{itemize}
  \item DeepSeek-R1 技术报告——RL 后训练的新范式
  \item OpenAI ``From Compute-Bound to Data-Bound'' Blog Post
  \item Visual Jigsaw / Spatial Intelligence Benchmark 相关工作
  \item Claude Code 产品设计与 Agentic 架构
  \item Slow Weight / Fast Weight 经典文献（Hinton et al.）
\end{itemize}

\end{document}
